\documentclass[twocolumn,10pt]{article}

%% ---- Packages -------------------------------------------------------
\usepackage[T1]{fontenc}
\usepackage[utf8]{inputenc}
\usepackage{amsmath,amssymb,amsthm}
\usepackage{mathtools}
\usepackage[margin=1in]{geometry}
\usepackage{abstract}
\usepackage{titlesec}
\usepackage{parskip}
\usepackage{booktabs}
\usepackage{graphicx}
\usepackage{xcolor}
\usepackage{hyperref}
\usepackage[numbers,sort&compress]{natbib}
\usepackage{fancyhdr}
\usepackage{etoolbox}
\usepackage{tikz-cd}

%% ---- Display math spacing -------------------------------------------
\setlength{\abovedisplayskip}{5pt}
\setlength{\belowdisplayskip}{5pt}
\setlength{\abovedisplayshortskip}{3pt}
\setlength{\belowdisplayshortskip}{3pt}

%% ---- Hyperref -------------------------------------------------------
\hypersetup{colorlinks=true,linkcolor=black,citecolor=black,urlcolor=black}

%% ---- Theorem environments -------------------------------------------
\newtheoremstyle{thmstyle}{6pt}{4pt}{\itshape}{}{\bfseries}{.}{ }{}
\newtheoremstyle{defstyle}{6pt}{4pt}{\normalfont}{}{\bfseries}{.}{ }{}

\theoremstyle{thmstyle}
\newtheorem{theorem}{Theorem}[section]
\newtheorem{corollary}[theorem]{Corollary}
\newtheorem{lemma}[theorem]{Lemma}
\newtheorem{proposition}[theorem]{Proposition}

\theoremstyle{defstyle}
\newtheorem{definition}[theorem]{Definition}
\newtheorem{remark}[theorem]{Remark}
\newtheorem{openp}[theorem]{Open Problem}

%% ---- Section formatting ---------------------------------------------
\titleformat{\section}{\normalfont\large\bfseries}{\thesection.}{0.5em}{}
\titleformat{\subsection}{\normalfont\normalsize\bfseries}{\thesubsection.}{0.5em}{}

%% ---- Header / footer ------------------------------------------------
\pagestyle{fancy}
\fancyhf{}
\fancyhead[L]{\small The Geometry of What Is Learned}
\fancyhead[R]{\small\thepage}
\renewcommand{\headrulewidth}{0.3pt}

%% ---- Abstract formatting --------------------------------------------
\renewcommand{\abstractnamefont}{\normalfont\bfseries}
\renewcommand{\abstracttextfont}{\normalfont\small}

%% ---- Custom commands ------------------------------------------------
\newcommand{\Xspace}{\mathcal{X}}
\newcommand{\Fprior}{\mathcal{F}}
\newcommand{\leff}{\lambda_{\mathrm{eff}}}
\newcommand{\Vdomain}{\mathcal{V}}
\newcommand{\Aset}{\mathcal{A}}

%% ====================================================================
\begin{document}

\twocolumn[{%
\begin{@twocolumnfalse}
\begin{center}
  {\LARGE\bfseries The Geometry of What Is Learned:\\
  Realization Maps, Fiber Structure, and the Semantic Manifold}\\[0.8em]
  {\large Flyxion}\\[0.2em]
  {\normalsize Independent Researcher}\\[0.2em]
  {\normalsize March 2026}
\end{center}
\begin{abstract}
The parameters of a neural network are not its meaning. This essay
develops the geometric consequences of that observation. The central
object is the realization map $\Phi : \Theta \to \mathcal{N}$, sending
parameter vectors to the functions they compute. Its fibers are the
equivalence classes of implementations of the same function, forming
smooth submanifolds of $\Theta$ whose dimension is determined by the
architectural symmetry group $G_\mathcal{A}$. The quotient
$\mathcal{N} \cong \Theta/G_\mathcal{A}$ is the semantic manifold: the
space of distinct realizable functions equipped with a canonical
Riemannian metric derived from the data distribution. Standard
quantities used to analyze learning---flatness, parameter magnitude,
gradient norms---are not invariant under $G_\mathcal{A}$ and carry no
semantic content. The correct geometry is that of $\mathcal{N}$:
learning is gradient flow on $(\mathcal{N}, g)$, generalization depends
on the semantic dimension $d = \dim\mathcal{N}$ and the connection
curvature $\kappa$, and compression is lossless if and only if it
preserves the fiber. Entanglement---the failure of the horizontal
distribution to be integrable, measured by the connection curvature
$\Omega$---simultaneously degrades optimization, generalization, and
compression capacity. The semantic manifold carries a canonical
statistical structure, a universal categorical characterization, a
gauge interpretation of redundancy, and a global topology that
constrains optimization. The correct unit of analysis is the equivalence
class under $\Phi$: meaning is invariant under implementation.
\end{abstract}
\vspace{1em}
\end{@twocolumnfalse}
}]

%% ====================================================================
\section{Introduction}
\label{sec:intro}

The parameters of a neural network are not its meaning. This
distinction, though obvious when stated, is almost never taken seriously
as a mathematical commitment. The prevailing practice treats the
parameter vector $\theta \in \Theta$ as the primary object of study:
loss functions are defined over $\Theta$, optimization trajectories are
traced through $\Theta$, generalization bounds are stated in terms of
properties of $\Theta$, and pruning methods operate by removing
coordinates of $\theta$. The implicit assumption is that the geometry of
$\Theta$ is the geometry that matters. It is not.

A neural network with parameters $\theta$ computes a function
$f_\theta : \mathcal{X} \to \mathcal{Y}$. Two parameter vectors $\theta$
and $\theta'$ that compute the same function are, from every behavioral
and semantic standpoint, identical. Yet the Euclidean geometry on
$\Theta$ treats them as distinct points, assigns them different gradient
norms, and would have an optimizer treat moving from one to the other as
meaningful. None of this is meaningful. It is an artifact of the
coordinate system, not a property of the model.

The failure runs deeper than a notational inconvenience. The standard
interpretation of flat minima as a signature of generalization is not
invariant under reparameterization of $\Theta$: a sharp minimum in one
parameterization can be made arbitrarily flat by a smooth coordinate
change that leaves $f_\theta$ unchanged. The same problem infects
pruning, where parameters with small magnitude are removed under the
assumption that they contribute little to the realized function---but
magnitude is coordinate-dependent and can be redistributed arbitrarily
within an equivalence class of identical functions.

The resolution begins with a single observation: there is a map
\[
\Phi : \Theta \to \mathcal{N}
\]
sending each parameter vector to the function it computes, where
$\mathcal{N}$ is the space of realizable functions. This map is the
fundamental object of study. Its fibers $\Phi^{-1}(f)$ are the sets of
all parameter vectors computing the same function. The quotient
$\mathcal{N} \cong \Theta/{\sim}$ is the space where meaning lives.

This essay develops the consequences of taking $\Phi$ seriously as a
mathematical object across its full geometric, statistical, categorical,
and topological dimensions. The argument proceeds through local
differential structure, global topology, statistical grounding,
categorical universality, variational derivation of learning, stability
theory, gauge interpretation of symmetry, and information geometry. The
goal throughout is not maximum generality but the minimal structure
sufficient to make the core claims precise.

%% ====================================================================
\section{System Model}
\label{sec:model}

\paragraph{Overview.}
This section collects all primitive objects into a single formal tuple
and introduces the notation used throughout.

\begin{table}[t]
\centering
\caption{Principal notation.}
\label{tab:notation}
\small
\begin{tabular}{@{}ll@{}}
\toprule
Symbol & Meaning \\
\midrule
$\Theta \subseteq \mathbb{R}^n$ & Parameter space \\
$\mathcal{X}, \mathcal{Y}$ & Input and output spaces \\
$\mathcal{F}(\mathcal{X},\mathcal{Y})$ & Function space \\
$\Phi : \Theta \to \mathcal{N}$ & Realization map \\
$\mathcal{N}$ & Semantic manifold \\
$G_\mathcal{A}$ & Architectural symmetry group \\
$\Phi^{-1}(f)$ & Implementation fiber over $f$ \\
$V_\theta = \ker D\Phi(\theta)$ & Vertical subspace \\
$H_\theta = V_\theta^\perp$ & Horizontal subspace \\
$g_f$ & Semantic metric on $\mathcal{N}$ \\
$G_\theta$ & Pullback metric on $\Theta$ \\
$\Omega$ & Connection curvature \\
$d = \dim\mathcal{N}$ & Semantic dimension \\
$\kappa = \|\Omega\|$ & Entanglement measure \\
$\mathcal{C}(f)$ & Policy complexity \\
$\mathcal{S}(f)$ & Policy entropy \\
$\mathcal{F}(f)$ & Semantic free energy \\
\bottomrule
\end{tabular}
\end{table}

\begin{definition}[Realization Map]
The \emph{realization map} of architecture $\mathcal{A}$ is
$\Phi : \Theta \to \mathcal{F}(\mathcal{X},\mathcal{Y})$, defined by
$\Phi(\theta) = f_\theta$. The \emph{semantic space}
$\mathcal{N} := \mathrm{image}(\Phi)$ is the set of all realizable
functions.
\end{definition}

\begin{definition}[Implementation Fiber]
For $f \in \mathcal{N}$, the \emph{implementation fiber} is
$\Phi^{-1}(f) = \{\theta \in \Theta \mid \Phi(\theta) = f\}$.
The equivalence relation $\theta \sim \theta'$ if and only if
$\Phi(\theta) = \Phi(\theta')$ partitions $\Theta$ and yields the
quotient $\mathcal{N} \cong \Theta/{\sim}$.
\end{definition}

\begin{definition}[Architectural Symmetry Group]
The \emph{symmetry group} $G_\mathcal{A}$ is the group of
diffeomorphisms $\sigma : \Theta \to \Theta$ satisfying
$\Phi \circ \sigma = \Phi$.
\end{definition}

\begin{proposition}[Smooth Quotient]
When $G_\mathcal{A}$ acts freely and properly on the regular locus
$\Theta^{\mathrm{reg}}$, the quotient $\mathcal{N}^{\mathrm{reg}} =
\Theta^{\mathrm{reg}}/G_\mathcal{A}$ is a smooth manifold of dimension
$n - \dim G_\mathcal{A}$, and $\Phi$ is a smooth principal
$G_\mathcal{A}$-bundle.
\end{proposition}

\begin{definition}[Horizontal and Vertical Subspaces]
At each regular $\theta$, define the \emph{vertical subspace}
$V_\theta := \ker D\Phi(\theta)$ and the \emph{horizontal subspace}
$H_\theta := V_\theta^\perp$, so that
$T_\theta\Theta = H_\theta \oplus V_\theta$.
\end{definition}

\begin{remark}[Geometric levels]
Four levels of structure are active simultaneously.
Level~(A): parameter space $\Theta \subset \mathbb{R}^n$,
finite-dimensional and Euclidean.
Level~(B): function space $\mathcal{F}(\mathcal{X},\mathcal{Y})$,
infinite-dimensional; we do not work here directly.
Level~(C): graph of $\Phi$, diffeomorphic to $\Theta$.
Level~(D): semantic space $\mathcal{N} = \mathrm{image}(\Phi)$, the
primary object of study. All claims about manifold structure and
learning dynamics refer to level~(D) unless stated otherwise.
\end{remark}

%% ====================================================================
\section{Differential Structure of the Realization Map}
\label{sec:differential}

\paragraph{Overview.}
All subsequent structure depends on the local behavior of $D\Phi(\theta)$.
This section establishes the rank stratification, introduces the
singular set, and derives a local normal form making the fiber geometry
explicit.

\begin{definition}[Rank and Regularity]
The \emph{rank} of $\Phi$ at $\theta$ is
$\mathrm{rank}(D\Phi(\theta)) := \dim\,\mathrm{image}(D\Phi(\theta))$.
A point $\theta$ is \emph{regular} if this rank equals $d = \dim\mathcal{N}$;
otherwise it is \emph{singular}. The \emph{singular set} is
$\Sigma := \{\theta \in \Theta \mid \mathrm{rank}(D\Phi(\theta)) < d\}$.
\end{definition}

\begin{proposition}[Generic Regularity]
Under mild smoothness assumptions on $\mathcal{A}$, the singular set
$\Sigma$ has measure zero in $\Theta$. This follows from Sard's theorem
applied to the smooth map $\Phi$, which for typical architectures is
polynomial or piecewise smooth.
\end{proposition}

All subsequent statements are understood to hold on the regular locus
$\Theta^{\mathrm{reg}} := \Theta \setminus \Sigma$.

\begin{proposition}[Local Normal Form]
\label{prop:local-normal-form}
Let $\theta \in \Theta^{\mathrm{reg}}$. Then there exist local
coordinates $(u,v)$ on a neighborhood $U \subset \Theta$ with
$u \in \mathbb{R}^d$, $v \in \mathbb{R}^{n-d}$, such that
$\Phi(u,v) = u$.
\end{proposition}

\begin{proof}[Sketch]
Since $\theta$ is regular, $D\Phi(\theta)$ has full rank $d$. By the
constant rank theorem, there exist coordinates in which $\Phi$ reduces
to projection onto the first $d$ coordinates.
\end{proof}

Three corollaries follow immediately.

\begin{corollary}[Explicit Fiber, Vertical Space, Semantic Dimension]
In the coordinates of Proposition~\ref{prop:local-normal-form}: the
fiber through $(u_0, v_0)$ is $\{(u_0, v) \mid v \in \mathbb{R}^{n-d}\}$;
the vertical space is $V_\theta = \mathrm{span}\{\partial/\partial v_i\}$
and the horizontal space is $H_\theta = \mathrm{span}\{\partial/\partial u_j\}$;
and the semantic dimension $d = \mathrm{rank}(D\Phi(\theta))$ is constant
on $\Theta^{\mathrm{reg}}$.
\end{corollary}

\begin{proposition}[Rank Stratification]
The parameter space decomposes into strata
$\Theta = \bigsqcup_{k=0}^d \Theta_k$ where
$\Theta_k := \{\theta \mid \mathrm{rank}(D\Phi(\theta)) = k\}$,
each a submanifold. Points in lower-rank strata correspond to degenerate
parameterizations: dead neurons, redundant layers, or collapsed
representations. The regular locus $\Theta^{\mathrm{reg}} = \Theta_d$
is open and dense.
\end{proposition}

The realization map behaves locally as a projection:
$\Theta \approx \mathbb{R}^d \times \mathbb{R}^{n-d} \xrightarrow{\Phi}
\mathbb{R}^d$. The first factor carries semantic degrees of freedom; the
second carries implementation degrees of freedom. All subsequent
constructions are global consequences of this local structure. At
singular points, the semantic dimension collapses, the fiber dimension
increases, and the geometry of $\mathcal{N}$ becomes non-smooth.

%% ====================================================================
\section{Fiber Geometry}
\label{sec:fiber}

\paragraph{Overview.}
Fibers are the geometric expression of implementation redundancy. Their
dimension, curvature, and relationship to the symmetry group determine
what optimization can do without changing the model.

At a regular $\theta$, the fiber has dimension
$n - \mathrm{rank}(D\Phi(\theta))$, which quantifies the irreducible
redundancy of the parameterization.

\begin{proposition}[Fiber-Orbit Coincidence]
If $G_\mathcal{A}$ acts freely and transitively on the fibers of $\Phi$
at regular points, then $\Phi^{-1}(\Phi(\theta)) = G_\mathcal{A}\cdot\theta$
for all regular $\theta$, and $\Theta \to \mathcal{N}$ is a principal
$G_\mathcal{A}$-bundle.
\end{proposition}

\begin{definition}[Fiber Curvature]
The \emph{fiber curvature} at $\theta$ is the second fundamental form
$II_\theta : V_\theta \times V_\theta \to H_\theta$ of the fiber
$\Phi^{-1}(\Phi(\theta))$ as a submanifold of $(\Theta, g)$.
It measures the rate at which the fiber's tangent space rotates as one
moves along it.
\end{definition}

\begin{proposition}[Pruning as Fiber Projection]
\label{prop:pruning-fiber}
A pruning map $P : \Theta \to \Theta' \subseteq \Theta$ is
\emph{lossless at $\theta$} if and only if $\Phi(P(\theta)) = \Phi(\theta)$,
i.e., $P$ moves $\theta$ within its fiber.
\end{proposition}

For ReLU networks, the activation pattern is locally constant within
each polyhedral region of $\Theta$, making $\Phi$ locally linear and
the fiber locally affine. Fiber curvature is concentrated at activation
boundaries, where the local linear structure of $\Phi$ changes
discontinuously.

%% ====================================================================
\section{The Failure of Parameter Geometry}
\label{sec:failure}

\paragraph{Overview.}
A quantity $Q : \Theta \to \mathbb{R}$ is a \emph{semantic invariant}
if $Q(\theta) = Q(\sigma(\theta))$ for all $\sigma \in G_\mathcal{A}$,
equivalently if $Q$ factors through $\Phi$. Three standard proxies for
semantic properties fail this test.

\begin{proposition}[Flatness is not a semantic invariant]
For any architecture with scaling symmetries, the spectrum of
$\nabla^2(\mathcal{L}\circ\Phi)(\theta^*)$ is not invariant under
$G_\mathcal{A}$. For any $M > 0$ there exists $\sigma \in G_\mathcal{A}$
with $\Phi(\sigma(\theta^*)) = \Phi(\theta^*)$ and
$\lambda_{\max}(\nabla^2(\mathcal{L}\circ\Phi)(\sigma(\theta^*))) > M$.
The flatness measure is therefore a property of the coordinate
representation, not of the function.
\end{proposition}

\begin{proof}[Sketch]
For a two-layer ReLU network, the transformation
$(w, v) \mapsto (\alpha w, v/\alpha)$ is in $G_\mathcal{A}$ for any
$\alpha > 0$ and leaves $f_\theta$ unchanged. The dominant term of the
Hessian scales as $\alpha^2$ in the rescaled neuron direction. Taking
$\alpha \to \infty$ drives $\lambda_{\max}$ to infinity while
$\Phi(\sigma_\alpha(\theta^*))$ remains fixed.
\end{proof}

\begin{proposition}[Parameter magnitude is not a semantic invariant]
For any architecture with a scaling symmetry and any $\epsilon > 0$,
there exists $\sigma \in G_\mathcal{A}$ with $\Phi(\sigma(\theta)) =
\Phi(\theta)$ and $\|\sigma(\theta)\|_1 < \epsilon$. Parameter magnitude
can be redistributed arbitrarily within a fiber.
\end{proposition}

\begin{proposition}[Gradient descent is not reparameterization-invariant]
Let $\psi : \Theta' \to \Theta$ be a smooth diffeomorphism that is not
an isometry. Then the gradient descent trajectory in $\Theta'$ and the
image under $\psi^{-1}$ of the trajectory in $\Theta$ are generically
distinct paths, even though they converge to the same semantic
minimizers.
\end{proposition}

\begin{proposition}[Non-invariance implies semantic vacuity]
\label{prop:vacuity}
Let $Q : \Theta \to \mathbb{R}$ not be constant on fibers of $\Phi$.
Then there exist $\theta_1, \theta_2$ with $\Phi(\theta_1) =
\Phi(\theta_2)$ and $Q(\theta_1) \neq Q(\theta_2)$. Any conclusion
drawn from $Q$ about $f = \Phi(\theta)$ is either coincidental or
assumes a canonical fiber representative.
\end{proposition}

Parameter space has two legitimate roles: as a computational substrate
for optimization, and as the total space of a fiber bundle over
$\mathcal{N}$. In both roles it is studied relative to $\Phi$, not as
an intrinsic Euclidean space.

%% ====================================================================
\section{The Quotient Manifold}
\label{sec:quotient}

\paragraph{Overview.}
This section establishes the smooth structure, metric, geodesics,
curvature, and free energy of $\mathcal{N}$.

\begin{definition}[Semantic Metric]
\label{def:semantic-metric}
Let $\mathcal{D}$ be a distribution on $\mathcal{X}$. For
$\delta f_1, \delta f_2 \in T_f\mathcal{N}$, define
\[
g_f(\delta f_1,\delta f_2) :=
\mathbb{E}_{x\sim\mathcal{D}}[\delta f_1(x)\,\delta f_2(x)].
\]
The \emph{semantic distance} is
$d(f_1,f_2) := (\mathbb{E}_{x\sim\mathcal{D}}[(f_1(x)-f_2(x))^2])^{1/2}$.
\end{definition}

\begin{proposition}[Pullback Metric and Fisher Information]
The pullback $G_\theta(u,v) := g_{\Phi(\theta)}(D\Phi(\theta)u,
D\Phi(\theta)v)$ equals the Fisher information matrix up to a positive
constant for probabilistic models \citep{amari1998natural}.
\end{proposition}

\begin{proposition}[Vertical Directions are Metrically Null]
For all $\theta \in \Theta^{\mathrm{reg}}$ and all $v \in V_\theta$,
$G_\theta(v,w) = 0$ for all $w$. The pullback metric descends to a
non-degenerate metric on $H_\theta \cong T_{\Phi(\theta)}\mathcal{N}$.
\end{proposition}

\begin{proposition}[Curvature Decomposition]
\label{prop:curvature-decomp}
Let $\tilde e_1, \tilde e_2$ be horizontal lifts of orthonormal
$e_1, e_2 \in T_f\mathcal{N}$. Then
\[
K^{\mathcal{N}}(e_1,e_2) = K^\Theta(\tilde e_1,\tilde e_2)
+ \|\Omega(\tilde e_1,\tilde e_2)\|^2 - B(\tilde e_1,\tilde e_2),
\]
where $\Omega$ is the connection curvature and $B$ depends on the
second fundamental form of the fibers
\citep{oneill1966fundamental}. The term $\|\Omega\|^2$ contributes
positively: fiber twisting increases the intrinsic curvature of
$\mathcal{N}$.
\end{proposition}

\begin{definition}[Semantic Free Energy]
\label{def:free-energy}
The \emph{policy complexity} is
$\mathcal{C}(f) := \int_\mathcal{X}\|\nabla_x f(x)\|^2\,d\mu(x)$,
the \emph{policy entropy} is
$\mathcal{S}(f) := -\int p_f(y)\log p_f(y)\,dy$, and the
\emph{semantic free energy} is
$\mathcal{F}(f) := \mathcal{C}(f) - \tau\mathcal{S}(f)$ for $\tau > 0$.
Low free energy indicates a favorable balance between simplicity and
informativeness.
\end{definition}

%% ====================================================================
\section{Pushforward Measures and Statistical Structure}
\label{sec:pushforward}

\paragraph{Overview.}
The semantic metric arises naturally from the pushforward of the data
distribution through $\Phi$. This section shows that the Fisher metric
is not an additional assumption but a consequence of the realization map,
and connects $\mathcal{N}$ to information geometry and optimal transport.

\begin{definition}[Model Distribution]
For $\theta \in \Theta$, define the \emph{model distribution}
$\mu_\theta := f_\theta{}_\#\mathcal{D}$, the pushforward of
$\mathcal{D}$ through $f_\theta$. The \emph{statistical realization map}
$\Phi_{\mathrm{stat}} : \Theta \to \mathcal{P}(\mathcal{Y})$,
$\Phi_{\mathrm{stat}}(\theta) = \mu_\theta$, factors through $\Phi$:
$\mu_\theta = \mu_{\theta'}$ whenever $\Phi(\theta) = \Phi(\theta')$.
Thus statistical structure depends only on the semantic class.
\end{definition}

\begin{proposition}[Second-Order Expansion of Divergence]
\label{prop:second-order}
Let $f_t$ be a smooth curve in $\mathcal{N}$ with $f_0 = f$. Then
\[
D(\mu_{f_t} \| \mu_f)
= \tfrac{1}{2} t^2\, g_f(\dot f,\dot f) + o(t^2),
\]
where $g_f$ is the semantic metric of Definition~\ref{def:semantic-metric}.
The semantic metric is therefore the infinitesimal form of statistical
distinguishability.
\end{proposition}

\begin{proof}[Sketch]
Expand the divergence to second order. The first variation vanishes at
$t = 0$, and the second variation yields a quadratic form that, by
direct computation, equals the Fisher information inner product. The
identification follows from Definition~\ref{def:semantic-metric}.
\end{proof}

\begin{proposition}[Metric Equivalence]
For probabilistic models with likelihood $p_\theta(y \mid x)$, the
pullback metric $G_\theta$ coincides up to a positive constant with the
Fisher information matrix
$\mathbb{E}_{(x,y)}[\nabla_\theta\log p_\theta(y\mid x)\,
\nabla_\theta\log p_\theta(y\mid x)^\top]$.
\end{proposition}

\begin{proposition}[Semantic Divergence is Well-Defined]
The divergence $D(f_1,f_2) := D(\mu_{f_1}\|\mu_{f_2})$ depends only on
$[f_1],[f_2] \in \mathcal{N}$ and is invariant under $G_\mathcal{A}$.
All statistical distinguishability lives on $\mathcal{N}$, not on
$\Theta$.
\end{proposition}

The Wasserstein-2 metric $W_2(\mu_{f_1},\mu_{f_2})$ defines an
alternative geometry on $\mathcal{N}$ based on optimal transport
\citep{villani2009optimal}, measuring global transport cost rather than
local distinguishability. Both are induced by $\Phi$ and are therefore
semantic. The geometry of $\mathcal{N}$ is simultaneously geometric,
functional, and statistical.

%% ====================================================================
\section{The Quotient as a Universal Construction}
\label{sec:categorical}

\paragraph{Overview.}
The semantic manifold $\mathcal{N}$ is characterized by a universal
property: it is the unique object through which all semantic invariants
factor. This places the invariance thesis on categorical footing.

The equivalence relation induced by $\Phi$ can be expressed as a fiber
product
$\Theta \times_{\mathcal{N}} \Theta :=
\{(\theta_1,\theta_2) \in \Theta \times \Theta \mid
\Phi(\theta_1) = \Phi(\theta_2)\}$,
with two canonical projections $\pi_1, \pi_2 :
\Theta \times_{\mathcal{N}} \Theta \rightrightarrows \Theta$.

\begin{proposition}[Coequalizer Characterization]
\label{prop:coequalizer}
The semantic space $\mathcal{N}$ is the coequalizer of the diagram
$\Theta \times_{\mathcal{N}} \Theta
\rightrightarrows \Theta \xrightarrow{\Phi} \mathcal{N}$.
For any space $Z$ and map $Q : \Theta \to Z$ satisfying
$Q \circ \pi_1 = Q \circ \pi_2$, there exists a unique map
$\tilde Q : \mathcal{N} \to Z$ with $Q = \tilde Q \circ \Phi$.
\end{proposition}

\begin{corollary}[Factorization of Invariants]
\label{cor:factorization}
Every semantic invariant $Q$ factors uniquely through $\mathcal{N}$.
The semantic manifold is the minimal representation of all semantically
meaningful quantities.
\end{corollary}

\begin{proposition}[Non-Invariance as Obstruction]
Let $Q : \Theta \to Z$ not be constant on fibers. Then there does not
exist any map $\tilde Q : \mathcal{N} \to Z$ with $Q = \tilde Q \circ \Phi$.
Non-invariant quantities cannot be interpreted as functions of the
learned model; they are functions of the implementation.
\end{proposition}

The following statements from earlier sections become immediate
corollaries of the universal property: Proposition~\ref{prop:vacuity}
holds because any non-invariant $Q$ fails to factor through $\mathcal{N}$;
the semantic metric $g$ is well-defined because it depends only on
$\Phi(\theta)$; and the pushforward measure $\mu_\theta$ depends only on
the equivalence class $[\theta]$. The quotient $\mathcal{N}$ is not a
convenient abstraction; it is the unique space in which semantically
meaningful quantities can be defined.

%% ====================================================================
\section{Learning as Geometric Motion}
\label{sec:learning}

\paragraph{Overview.}
This section analyzes what optimization does when viewed through $\Phi$,
showing that gradient descent decomposes into horizontal (semantic) and
vertical (redundant) components, and that the correct algorithm is
natural gradient descent.

\begin{proposition}[Gradient Descent is Horizontal]
The standard gradient
$\nabla_\theta(\mathcal{L}\circ\Phi)(\theta) =
D\Phi(\theta)^*\nabla_f\mathcal{L}$ lies entirely in
$H_\theta = (\ker D\Phi(\theta))^\perp$.
The vertical component vanishes identically.
\end{proposition}

\begin{proof}
The image of $D\Phi(\theta)^*$ is $(\ker D\Phi(\theta))^\perp = H_\theta$,
since $\mathrm{im}(A^*) = (\ker A)^\perp$ for any linear map between
inner product spaces.
\end{proof}

Despite this, standard gradient descent is not semantically correct:
the horizontal component corresponds to a direction in $T_f\mathcal{N}$
via $D\Phi(\theta)$, but when the horizontal subspace is defined by the
Euclidean metric on $\Theta$ rather than the semantic metric $g$, the
resulting motion in $\mathcal{N}$ is not the Riemannian gradient of
$\mathcal{L}$. This is the core mismatch.

\begin{definition}[Natural Gradient]
The \emph{natural gradient} update is
$\dot\theta = -G_\theta^+\nabla_\theta\ell(\theta)$,
where $G_\theta^+ = D\Phi(\theta)^+(g^{-1}_{\Phi(\theta)})D\Phi(\theta)^{+*}$
is the pseudoinverse of the pullback metric.
\end{definition}

\begin{proposition}[Natural Gradient is Semantically Correct]
The natural gradient satisfies: $\dot\theta \in H_\theta$;
$D\Phi(\theta)\dot\theta = -\mathrm{grad}_g\mathcal{L}(\Phi(\theta))$;
and the update is $G_\mathcal{A}$-invariant. Standard gradient descent
is not $G_\mathcal{A}$-invariant.
\end{proposition}

\begin{proposition}[Vertical Drift and Connection Curvature]
For a gradient step of size $\eta$,
\[
(\theta' - \theta)_V =
-\tfrac{\eta^2}{2}\,\Omega(\dot\theta,\dot\theta) + O(\eta^3).
\]
High connection curvature generates significant vertical drift per step,
causing the optimizer to wander within fibers while the semantic
trajectory converges.
\end{proposition}

\begin{proposition}[Curvature and Effective Learning Rate]
In a region with $|K^\mathcal{N}| \leq K_0$, the effective step size
for guaranteed semantic progress satisfies
$\eta_{\mathrm{eff}} \leq (L + K_0\,d(f,f^*))^{-1}$.
High curvature forces smaller steps and degrades the convergence rate.
\end{proposition}

%% ====================================================================
\section{Projected Dynamics and Variational Principle}
\label{sec:projected}

\paragraph{Overview.}
This section derives the natural gradient from a constrained variational
principle, removing any arbitrariness from the learning rule and
establishing that learning is a projection problem induced by $\Phi$.

The desired semantic evolution is $\dot f = -\mathrm{grad}_g\mathcal{L}(f)$,
which induces the constraint on parameter velocities:
$D\Phi(\theta)\,\dot\theta = -\mathrm{grad}_g\mathcal{L}(\Phi(\theta))$.
Among all velocities satisfying this constraint, the natural choice is
the one of minimal norm.

\begin{proposition}[Constrained Variational Characterization]
\label{prop:variational}
The natural gradient flow is the solution to
\[
\min_{\dot\theta \in T_\theta\Theta} \|\dot\theta\|^2
\quad \text{subject to} \quad
D\Phi(\theta)\,\dot\theta = -\mathrm{grad}_g\mathcal{L}(\Phi(\theta)).
\]
\end{proposition}

\begin{proof}
Introduce a Lagrange multiplier $\lambda \in T_f^*\mathcal{N}$ and
consider $\mathcal{J}(\dot\theta,\lambda) = \|\dot\theta\|^2 +
\langle\lambda, D\Phi(\theta)\dot\theta + \mathrm{grad}_g\mathcal{L}\rangle$.
Stationarity in $\dot\theta$ gives
$2\dot\theta + D\Phi(\theta)^*\lambda = 0$, so
$\dot\theta = -\tfrac{1}{2}D\Phi(\theta)^*\lambda$.
Substituting into the constraint gives
$D\Phi(\theta)D\Phi(\theta)^*\lambda = 2\,\mathrm{grad}_g\mathcal{L}$,
determining $\lambda$. Substituting back yields the natural gradient.
\end{proof}

\begin{corollary}[Natural Gradient as Pseudoinverse Projection]
The solution is $\dot\theta = -D\Phi(\theta)^+\mathrm{grad}_g\mathcal{L}$,
where $D\Phi(\theta)^+ := D\Phi(\theta)^*(D\Phi(\theta)D\Phi(\theta)^*)^{-1}$
is the Moore-Penrose pseudoinverse. The natural gradient is the
least-squares solution to the semantic descent equation.
\end{corollary}

\begin{proposition}[Projection Interpretation]
The operator $\Pi_\theta := D\Phi(\theta)^*(D\Phi(\theta)D\Phi(\theta)^*)^{-1}
D\Phi(\theta)$ is the orthogonal projector onto $H_\theta$. Learning
proceeds by projecting the gradient onto $H_\theta$ and discarding all
vertical components.
\end{proposition}

\begin{proposition}[Uniqueness of the Natural Gradient]
Among all parameter updates satisfying
$D\Phi(\theta)\dot\theta = -\mathrm{grad}_g\mathcal{L}$, the natural
gradient is the unique update minimizing $\|\dot\theta\|$. It is
therefore not a heuristic but a canonical construction.
\end{proposition}

\begin{proposition}[First-Order Consistency]
For step size $\eta$, the discrete update satisfies
$\Phi(\theta_{k+1}) = \Phi(\theta_k) - \eta\,\mathrm{grad}_g\mathcal{L}
+ O(\eta^2)$, inducing the correct first-order semantic motion.
\end{proposition}

Standard gradient descent fails because it does not solve the
constrained variational problem: it uses the Euclidean metric to measure
parameter velocity instead of finding the minimum-norm velocity that
achieves the desired semantic change. The geometry of $\Phi$ determines
both the constraint and its solution.

%% ====================================================================
\section{Identifiability and Observability}
\label{sec:identifiability}

\paragraph{Overview.}
The realization map determines not only semantic content but everything
that can be observed. The equivalence relation induced by $\Phi$
coincides with observational equivalence: no experiment on inputs and
outputs can distinguish points within the same fiber.

\begin{definition}[Observational Equivalence]
Parameters $\theta, \theta' \in \Theta$ are \emph{observationally
equivalent} if $f_\theta(x) = f_{\theta'}(x)$ for all $x \in \mathcal{X}$.
\end{definition}

\begin{proposition}[Equivalence with Fiber Relation]
$\theta \sim_{\mathrm{obs}} \theta'$ if and only if
$\Phi(\theta) = \Phi(\theta')$.
\end{proposition}

Fibers are exactly the sets of parameters indistinguishable by any
input-output experiment.

\begin{proposition}[Failure of Identifiability]
For architectures with nontrivial $G_\mathcal{A}$, statistical
identifiability fails: $\sigma(\theta) \neq \theta$ but
$\mu_{\sigma(\theta)} = \mu_\theta$ for all $\sigma \in G_\mathcal{A}$.
The model is identifiable only up to equivalence classes, and this is
the strongest identifiability condition that can hold.
\end{proposition}

\begin{theorem}[Indistinguishability under $\Phi$]
\label{thm:indistinguishability}
Let $O$ be any observer with access only to input-output pairs
$(x, f_\theta(x))$ drawn from $\mathcal{D}$. For any $\theta, \theta'$
with $\Phi(\theta) = \Phi(\theta')$, the distributions of observations
are identical: $(x,f_\theta(x)) \overset{d}{=} (x,f_{\theta'}(x))$.
No statistical test can distinguish $\theta$ from $\theta'$,
regardless of dataset size, test procedure, or computational resources.
\end{theorem}

\begin{proof}
$\Phi(\theta) = \Phi(\theta')$ implies $f_\theta = f_{\theta'}$, so the
joint distribution of $(x,f_\theta(x))$ is identical.
\end{proof}

\begin{proposition}[Observables are Semantic Invariants]
Every function $Q : \Theta \to \mathbb{R}$ estimable from input-output
data is constant on fibers and therefore factors through $\Phi$.
Observable and semantic invariant are equivalent conditions.
\end{proposition}

The invariance thesis thus has an observational reading: meaning is not
what remains after ignoring implementation---it is what remains after
exhausting all possible observations. Parameter norms, individual
weights, and neuron activations tied to specific parameterizations are
unobservable and cannot be inferred from data.

%% ====================================================================
\section{Stability and Perturbation Theory}
\label{sec:stability}

\paragraph{Overview.}
The geometry induced by $\Phi$ determines how sensitive model outputs
are to perturbations. This section develops a perturbation theory for
$\Phi$, connecting stability to the singular structure of $D\Phi$ and
the curvature of $\mathcal{N}$.

For a perturbation $\delta\theta$, the induced change in function space
is $\delta f = D\Phi(\theta)\,\delta\theta$. A perturbation $\delta\theta$
has zero semantic effect if and only if $\delta\theta \in V_\theta$;
vertical directions are exactly the zero-sensitivity directions.

\begin{proposition}[Sensitivity Spectrum]
Let $D\Phi(\theta) = U\Sigma V^\top$ be the singular value
decomposition with singular values $\sigma_1 \geq \cdots \geq \sigma_d > 0$.
Large singular values correspond to highly sensitive directions; small
singular values to weakly sensitive ones; and zero singular values to
redundant (vertical) directions. The condition number
$\kappa(\theta) := \sigma_{\max}/\sigma_{\min}$ controls the anisotropy
of sensitivity and the stability of inversion.
\end{proposition}

\begin{proposition}[Second-Order Effects]
Even when $\delta\theta \in V_\theta$, higher-order terms may induce
nonzero semantic changes: $D\Phi(\theta)\delta\theta = 0$ does not
imply $\Phi(\theta + \delta\theta) = \Phi(\theta)$. The second-order
expansion is
$\Phi(\theta + \delta\theta) = \Phi(\theta) + D\Phi(\theta)\delta\theta
+ \tfrac{1}{2}D^2\Phi(\theta)(\delta\theta,\delta\theta) + O(\|\delta\theta\|^3)$,
so vertical directions are only infinitesimally null.
\end{proposition}

\begin{proposition}[Inverse Stability]
The minimal-norm solution to $D\Phi(\theta)\delta\theta = \delta f$ is
$\delta\theta = D\Phi(\theta)^+\delta f$, satisfying
$\|\delta\theta\| \leq \sigma_{\min}^{-1}\|\delta f\|$.
Small $\sigma_{\min}$ produces large parameter changes for small
functional changes.
\end{proposition}

\begin{proposition}[Noise Filtering and Geodesic Deviation]
If parameters are perturbed by noise $\delta\theta$, the output
variance is $\mathrm{Var}(\delta f) = D\Phi(\theta)\,\Sigma_\theta\,
D\Phi(\theta)^\top$. Vertical noise is completely filtered:
$\delta\theta \in V_\theta$ implies $\mathrm{Var}(\delta f) = 0$.
Curvature of $\mathcal{N}$ introduces geodesic deviation:
$D^2\delta f/dt^2 = -R(\dot\gamma,\delta f)\dot\gamma$, so positive
curvature contracts nearby trajectories and negative curvature causes
divergence.
\end{proposition}

The realization map acts as a filter: it removes redundant directions,
attenuates weakly sensitive ones, and amplifies sensitive ones.
Stability is therefore a property of how $\Phi$ transforms parameter
perturbations into functional changes, not a property of parameters per
se. Flat regions---where $\|D\Phi(\theta)\|$ is small in relevant
directions---correspond to low sensitivity, robustness to noise, and
improved generalization.

%% ====================================================================
\section{Gauge Symmetry and Redundancy}
\label{sec:gauge}

\paragraph{Overview.}
The redundancy of parameter space is structured by symmetry. This
section shows that the fibers of $\Phi$ arise from a gauge symmetry
and that learning dynamics have a natural gauge-fixed interpretation.

\begin{proposition}[Infinitesimal Gauge Generators are Vertical]
For $\xi \in \mathfrak{g} = \mathrm{Lie}(G_\mathcal{A})$, define the
vector field $X_\xi(\theta) = (d/dt)|_{t=0}\exp(t\xi)\cdot\theta$.
Then $X_\xi(\theta) \in V_\theta = \ker D\Phi(\theta)$.
\end{proposition}

\begin{proof}
Since $\Phi(\exp(t\xi)\cdot\theta) = \Phi(\theta)$ for all $t$,
differentiating at $t=0$ gives $D\Phi(\theta)X_\xi(\theta) = 0$.
\end{proof}

Vertical directions are therefore generated by infinitesimal symmetries.
Gauge equivalence implies semantic equivalence:
$\theta \sim_{\mathrm{gauge}} \theta'$ implies $\Phi(\theta) =
\Phi(\theta')$, though the converse need not hold if the full fiber is
larger than the group orbit.

\begin{definition}[Connection as Gauge Fixing]
A choice of horizontal subspace $H_\theta$ complementing $V_\theta$ is
a \emph{connection} on the bundle $\Theta \to \mathcal{N}$. Choosing
$H_\theta$ is equivalent to selecting a gauge: it picks a unique
representative of each equivalence class for infinitesimal motion.
Natural gradient descent is motion in a fixed gauge.
\end{definition}

\begin{proposition}[Curvature as Gauge Obstruction]
The curvature $\Omega(X,Y) = \mathrm{proj}_V([X^H,Y^H])$ is the
obstruction to integrability of the horizontal distribution. Nonzero
curvature implies that parallel transport is path-dependent and that
no global gauge fixing exists. This corresponds to entanglement between
semantic and implementation degrees of freedom.
\end{proposition}

\begin{proposition}[Gauge-Projected Dynamics]
The natural gradient flow satisfies $\dot\theta \in H_\theta$ at every
point: learning is constrained to remain orthogonal to all gauge
directions, eliminating gauge redundancy at the level of dynamics.
\end{proposition}

Every gauge-invariant function descends to $\mathcal{N}$, and gauge
invariance is sufficient for semantic invariance. The central thesis
acquires a gauge-theoretic reading: parameters are coordinates in a
gauge space; models are points in the quotient.

%% ====================================================================
\section{Generalization and Flat Minima}
\label{sec:generalization}

\paragraph{Overview.}
Generalization depends on the semantic dimension of the hypothesis class
and the connection curvature, not on the parameter count.

\begin{definition}[Semantic Hypothesis Class]
$\mathcal{H}_{B,\kappa} := \{f \in \mathcal{N} \mid \mathcal{F}(f)
\leq B,\; \|\Omega_f\| \leq \kappa\}$, where $\mathcal{F}$ is the
semantic free energy of Definition~\ref{def:free-energy}.
\end{definition}

\begin{proposition}[Semantic Dimension Reduction]
$\log N(\varepsilon,\mathcal{H}_{B,\kappa},d) \leq
C_1 d\log(A_B/\varepsilon) + C_2\kappa\,\varepsilon^{-1}$.
The covering number depends on the semantic dimension $d$, not the
parameter count $n$.
\end{proposition}

\begin{theorem}[Semantic Generalization Bound]
With probability $\geq 1-\delta$, uniformly for all
$f \in \mathcal{H}_{B,\kappa}$,
\[
\mathrm{GenErr}(f) \leq \inf_{\varepsilon>0}\!\left\{
\varepsilon + \sqrt{\frac{C_1 d\log(A_B/\varepsilon)
+ C_2\kappa\,\varepsilon^{-1} + \log(1/\delta)}{n}}
\right\}.
\]
\end{theorem}

When $\kappa$ is small the bound reduces to $\sqrt{d\log n/n}$
(dimension-dominant regime); when $\kappa$ is large the bound degrades
to $(\kappa/n)^{1/3}$ (curvature-dominant regime), showing that
entanglement worsens the generalization rate independently of model
size.

\begin{theorem}[Rademacher Bound on the Semantic Class]
Under the assumptions of the hypothesis class definition,
\[
\mathrm{GenErr}(f) \leq
\frac{C'L}{\sqrt{n}}\!\left(\sqrt{d}
+ \sqrt{\kappa}\,A_B^{1/2}\right)
+ \sqrt{\frac{\log(1/\delta)}{n}},
\]
up to logarithmic factors. The two terms measure semantic complexity
and geometric entanglement respectively.
\end{theorem}

\begin{definition}[Semantic Flatness]
A minimizer $f^* \in \mathcal{N}$ is \emph{semantically flat} if
$\nabla^2_g\mathcal{L}(f^*)$ has small eigenvalues in the metric $g$.
This is $G_\mathcal{A}$-invariant; the flat minima heuristic in $\Theta$
is its coordinate-dependent approximation, correct only when $D\Phi(\theta)$
is approximately an isometry.
\end{definition}

%% ====================================================================
\section{Projection and Compression}
\label{sec:compression}

\paragraph{Overview.}
The realization map induces a compression from parameter space to
semantic space. This section formalizes all compression schemes as
instances of a single construction and derives the geometric criterion
for lossless compression.

\begin{definition}[Compression Map]
A \emph{compression map} $\rho : \Theta \to \Theta'$ is
\emph{lossless at $\theta^*$} if $\Phi(\rho(\theta^*)) = \Phi(\theta^*)$.
The \emph{compression error} is
$\mathcal{E}(\rho,\theta^*) := d(\Phi(\rho(\theta^*)),\Phi(\theta^*))$.
\end{definition}

Pruning, distillation, low-rank factorization, and quantization are
all instances: they differ in the form of $\Theta'$ and the degree to
which they are fiber-aware. Distillation directly minimizes
$\mathcal{E}$; the others use coordinate-dependent proxies. A
compression is \emph{fiber-aware} if $\rho(\theta^*) - \theta^*$ lies
approximately in $V_{\theta^*}$; in that case
$\mathcal{E}(\rho,\theta^*) \leq \delta\|D\Phi(\theta^*)\|\cdot
\|\rho(\theta^*)-\theta^*\|$ for small $\delta > 0$.

\begin{proposition}[Lossless Compression Existence]
A lossless compression of $\theta^*$ into $\Theta'$ exists if and
only if $\Phi^{-1}(\Phi(\theta^*)) \cap \Theta' \neq \emptyset$.
\end{proposition}

\begin{proposition}[Universal Compression Criterion]
$\Phi$ factors through $\Theta'$ (universally lossless compression)
if and only if the fibers of $\rho$ contain the fibers of $\Phi$:
$\ker(\rho_*) \supseteq V_\theta$ at every regular point.
\end{proposition}

\begin{proposition}[Information Bottleneck]
The differential $D\Phi(\theta)$ acts as a local information bottleneck:
its rank is the effective dimension of information flow, and directions
with small singular values are attenuated. Fibers encode lost
information; large fibers correspond to high redundancy and improved
robustness.
\end{proposition}

\begin{proposition}[Fiber Size, Compression, and Generalization]
The fiber dimension $n - d$ simultaneously upper-bounds the lossless
compression gain and lower-bounds the generalization improvement from
working in $\mathcal{N}$ rather than $\Theta$. The geometry of the
fiber governs both phenomena through a single object.
\end{proposition}

%% ====================================================================
\section{Entanglement and Curvature}
\label{sec:entanglement}

\paragraph{Overview.}
This section examines the condition in which the fiber is geometrically
twisted so that semantic and implementation degrees of freedom cannot be
cleanly separated. Entanglement is measured by the connection curvature
$\Omega$ and the second fundamental form $II$.

\begin{proposition}[Entanglement as Non-integrability]
\label{prop:entanglement}
The horizontal distribution $H_\theta$ is integrable if and only if
$\Omega \equiv 0$. When $\Omega \neq 0$, horizontal vector fields fail
to close under the Lie bracket, so $\Theta$ does not locally decompose
as $\mathcal{N} \times G_\mathcal{A}$ and semantic motion necessarily
generates vertical displacement.
\end{proposition}

\begin{proof}[Sketch]
Frobenius' theorem: a distribution is integrable if and only if it is
involutive. The obstruction to involutivity for a principal bundle
connection is exactly the curvature $\Omega$.
\end{proof}

\begin{proposition}[Holonomy and Implementation Drift]
Let $\gamma$ be a loop in $\mathcal{N}$ based at $f$. The horizontal
lift starting at $\theta \in \Phi^{-1}(f)$ returns to
$\tilde\gamma(1) = \sigma_\gamma(\theta)$, where $\sigma_\gamma
\in G_\mathcal{A}$ is the holonomy element. Even when the semantic
state returns to its starting point, the implementation may not. The
fiber is dynamically active under learning.
\end{proposition}

\begin{proposition}[Curvature Obstructs Factorization]
If $\Omega \neq 0$ in a neighborhood of $\theta^*$, then no coordinate
system on $\Theta$ exists in which $\Phi$ decomposes locally as
$\Phi(\theta_{\mathrm{sem}}, \theta_{\mathrm{impl}}) =
\Phi(\theta_{\mathrm{sem}})$. Any factored approximation incurs
semantic error proportional to $\|\Omega\|$.
\end{proposition}

By the O'Neill formula (Proposition~\ref{prop:curvature-decomp}),
$\|\Omega\|^2$ contributes positively to $K^\mathcal{N}$: nearby
semantic directions couple under transport and geodesics diverge from
linear approximations more rapidly in entangled regions.

\begin{proposition}[Entanglement Degrades Optimization, Generalization, and Compression]
When $\|\Omega\|$ is large: horizontal updates induce vertical drift
of order $\eta^2\|\Omega\|$; the discrepancy operator becomes
ill-conditioned; gradient directions misalign with semantic gradients;
covering numbers and Rademacher complexity increase via the curvature
term $\sqrt{\kappa}A_B^{1/2}/\sqrt{n}$; and fiber navigation becomes
geometrically costly, preventing compression schemes from finding
optimal fiber representatives.
Two models with identical parameter counts and similar training loss
can differ dramatically in generalization: the difference is curvature,
not size.
\end{proposition}

\begin{definition}[Geometric Disentanglement]
A representation is \emph{disentangled} in a region if
$\|\Omega\| \approx 0$ and $\|II\| \approx 0$. In this regime fibers
are flat, horizontal motion is path-independent, compression is
fiber-aware and efficient, and optimization aligns with semantics.
Normalization schemes, residual connections, and prescribed
architectural symmetries function as curvature-reduction mechanisms
in precisely this sense.
\end{definition}

Entanglement is not an abstract property of representations. It is a
measurable geometric quantity---the curvature of the realization
bundle---that determines whether redundancy can be exploited, whether
learning aligns with meaning, and whether models can be simplified
without loss.

%% ====================================================================
\section{Global Topology and Nontrivial Structure of $\mathcal{N}$}
\label{sec:topology}

\paragraph{Overview.}
All preceding sections describe local properties of $\Phi$. This
section studies the global topology induced by $\Phi$, establishing
that $\mathcal{N}$ is a stratified space with nontrivial structure that
constrains optimization at a topological level.

The quotient topology on $\mathcal{N}$ is the finest topology making
$\Phi$ continuous. It may fail to be Hausdorff if fibers intersect
nontrivially under limits. The rank stratification of
Section~\ref{sec:differential} induces a stratification
$\mathcal{N} = \bigsqcup_{k=0}^d \mathcal{N}_k$ where
$\mathcal{N}_k := \Phi(\Theta_k)$ and each $\mathcal{N}_k$ is a
manifold of dimension $k$.

\begin{proposition}[Topological Obstruction to Optimization]
$\mathcal{N}$ may have multiple connected components, corresponding to
function classes that cannot be continuously deformed into one another.
If $f_1$ and $f_2$ lie in different components, no continuous
optimization path in $\Theta$ can transform one into the other without
leaving $\mathrm{image}(\Phi)$. Optimization is therefore constrained
by topology, not only by geometry.
\end{proposition}

Examples of distinct components include classifiers with different
decision boundary topology, functions with different global symmetry
properties, and representations with distinct topological invariants.

\begin{proposition}[Holonomy and Homotopy]
Parallel transport of $\theta$ along a loop $\gamma$ in $\mathcal{N}$
returns to $\sigma_\gamma(\theta) \in \Phi^{-1}(\Phi(\theta))$:
loops in $\mathcal{N}$ correspond to gauge transformations in $\Theta$.
Nontrivial homotopy classes in $\mathcal{N}$ indicate topological holes
or higher-dimensional features corresponding to global invariants of
function space.
\end{proposition}

\begin{proposition}[Morse-Theoretic Constraints]
The loss functional $\mathcal{L} : \mathcal{N} \to \mathbb{R}$ defines
a landscape over $\mathcal{N}$. The number and type of critical points
are constrained by Morse inequalities applied to the topology of
$\mathcal{N}$. The topological complexity of $\mathcal{N}$ therefore
imposes lower bounds on the number of critical points of any training
objective.
\end{proposition}

The semantic manifold $\mathcal{N}$ is not a simple smooth space. It is
stratified, possibly non-Hausdorff, and carries nontrivial homotopy.
Learning is navigation on a compressed, curved, and topologically
structured space of functions, constrained not only by local geometry
but by global invariants. This completes the picture: local structure
(Sections~\ref{sec:differential}--\ref{sec:stability}) and gauge
structure (Section~\ref{sec:gauge}) are embedded in a global topology
that governs what can be reached at all.

%% ====================================================================
\section{Meaning as Invariant}
\label{sec:meaning}

\paragraph{Overview.}
This section states what all preceding sections are collectively
establishing. The claim is philosophical in form but carries precise
mathematical content: meaning is invariant under implementation, and
the correct unit of analysis is the equivalence class under $\Phi$.

The invariance thesis has teeth because it is violated systematically in
current practice. Loss landscapes are analyzed in $\Theta$, not
$\mathcal{N}$. Generalization bounds are stated in terms of parameter
norms. Pruning is performed by magnitude thresholding. Interpretability
research analyzes individual neurons. Each produces the same error: a
coordinate-dependent quantity is mistaken for a semantic property.

\begin{definition}[Semantic Content]
The \emph{semantic content} of $\theta$ is the point $\Phi(\theta)
\in \mathcal{N}$. Two parameters have the same semantic content if and
only if they lie in the same fiber.
\end{definition}

\begin{theorem}[Realization Principle]
\label{thm:realization-principle}
All quantities invariant under $G_\mathcal{A}$ factor uniquely through
$\Phi$ and are functions on $\mathcal{N}$. Conversely, any non-invariant
quantity cannot be interpreted as a property of the learned function:
it is a property of the implementation.
\end{theorem}

\begin{proof}
The first claim is Corollary~\ref{cor:factorization}. The second is
Proposition~\ref{prop:vacuity} together with the universal property of
the coequalizer (Proposition~\ref{prop:coequalizer}).
\end{proof}

Working in $\mathcal{N}$ does not require choosing a canonical fiber
representative. It requires performing operations whose output is
independent of which representative is chosen. Natural gradient descent
satisfies this by construction. Distillation satisfies this by
minimizing semantic distance directly. The generalization bound satisfies
this by depending on $d$ and $\kappa$, both defined on $\mathcal{N}$.
The indistinguishability theorem (Theorem~\ref{thm:indistinguishability})
gives the observational reading: meaning is what remains after exhausting
all possible observations, not merely what remains after ignoring
implementation.

Everything in the essay compresses into a single chain:
\begin{equation}
\Phi(\theta_1) = \Phi(\theta_2)
\;\Rightarrow\; [\theta_1] = [\theta_2]
\;\Rightarrow\; \text{same meaning}
\;\Rightarrow\; \text{same generalization, compression, and dynamics.}
\label{eq:invariance-chain}
\end{equation}

%% ====================================================================
\section{Connections and Open Problems}
\label{sec:connections}

\paragraph{Overview.}
The framework connects to information geometry, algebraic geometry,
the loss landscape literature, equivariant architectures, mechanistic
interpretability, optimal transport, and gauge theory.

\begin{openp}[Information-geometric characterization of $\Omega$]
Express the connection curvature $\Omega$ in terms of the
$\alpha$-connection structure of information geometry
\citep{amari1998natural} and determine whether entanglement conditions
correspond to known properties of statistical model classes.
\end{openp}

\begin{openp}[Algebraic fiber structure]
For architectures with polynomial activations, compute the degree and
dimension of the fibers of $\Phi$ as algebraic varieties
\citep{kileel2019expressive} and extend to ReLU networks via
semialgebraic geometry.
\end{openp}

\begin{openp}[Semantic local minima]
Determine whether every local minimum of $\ell$ in $\Theta$ projects
to a local minimum of $\mathcal{L}$ in $\mathcal{N}$, whether spurious
local minima correspond to degenerate implementations, and whether
Morse theory on $\mathcal{N}$ accounts for the topology of sublevel
sets of $\ell$ on $\Theta$ \citep{bhojanapalli2016global}.
\end{openp}

\begin{openp}[Optimal symmetry group design]
Given target $\mathcal{N}^*$ and parameter budget $n$, characterize
symmetry groups $G$ achieving $\mathcal{N} \supseteq \mathcal{N}^*$
while maximizing $\dim G$ \citep{zaheer2017deep,thomas2018tensor}.
\end{openp}

\begin{openp}[Invariant feature construction]
Develop systematic methods for constructing $G_\mathcal{A}$-invariant
quantities from network weights and determine when such invariants
separate points of $\mathcal{N}$ \citep{olah2020zoom,elhage2021mathematical}.
\end{openp}

\begin{openp}[Wasserstein structure of $\mathcal{N}$]
For probabilistic models, determine the relationship between the semantic
metric $g$ and the Wasserstein-2 metric on model distributions, and
connect the semantic free energy $\mathcal{F}$ to displacement convexity
in Wasserstein space \citep{villani2009optimal}.
\end{openp}

\begin{openp}[Gauge-theoretic optimization]
Formulate natural gradient as a gauge-covariant gradient flow and
determine whether Yang-Mills equations have a meaningful analogue in the
optimization of $\mathcal{L}$ on $\mathcal{N}$.
\end{openp}

\begin{openp}[Global topology of $\mathcal{N}$]
Compute the fundamental group $\pi_1(\mathcal{N})$ and higher homotopy
groups for specific architectures. Determine whether non-trivial loops
correspond to observable training phenomena and whether the Morse
theory of $\mathcal{L}$ on $\mathcal{N}$ provides a complete account of
optimization landscape topology.
\end{openp}

\begin{openp}[Effective computation on $\mathcal{N}$]
Develop efficient algorithms for optimization, generalization assessment,
compression, and interpretation directly on $\mathcal{N}$ without
computing the full Jacobian $D\Phi(\theta)$. This is the primary
technical obstacle between the theoretical framework and its practical
application.
\end{openp}

%% ====================================================================
\section{Conclusion: The Geometry of What Is Learned}
\label{sec:conclusion}

The parameters of a neural network are not its meaning. From this
observation, pursued through differential geometry, probability theory,
category theory, variational analysis, gauge theory, and global
topology, a complete geometric theory has emerged. Every result in the
preceding sections is a consequence of taking the realization map
$\Phi : \Theta \to \mathcal{N}$ seriously as a mathematical object.

The fibers of $\Phi$ are smooth submanifolds of $\Theta$ with dimension
determined by $G_\mathcal{A}$. The local normal form makes fibers
explicit as coordinate planes, and Sard's theorem guarantees they are
well-behaved generically. The quotient $\mathcal{N} \cong
\Theta/G_\mathcal{A}$ carries the Fisher information metric,
characterized as the infinitesimal form of statistical distinguishability
via the pushforward of the data distribution. The universal property of
the coequalizer makes the invariance thesis categorically inevitable:
semantic invariants are precisely the quantities that factor through
$\Phi$, and non-invariants are provably unobservable.

Learning is the unique minimum-norm solution to a constrained variational
problem---natural gradient descent---derived without appeal to
heuristics. Stability is governed by the singular values of $D\Phi$ and
the curvature of $\mathcal{N}$, with vertical directions filtering noise
and curvature causing geodesic deviation. The symmetry group $G_\mathcal{A}$
generates the vertical directions as infinitesimal gauge transformations,
and the connection defines a gauge-fixed dynamics. Generalization depends
on $d = \dim\mathcal{N}$ and the connection curvature $\kappa$;
compression is lossless iff it preserves fibers; entanglement---measured
by $\Omega$---degrades all three simultaneously. The semantic manifold
carries a global topology that constrains what can be reached by any
continuous optimization.

The framework implies three things for practice. The relevant dimension
is $d = n - \dim G_\mathcal{A}$, not $n$. The symmetry group $G_\mathcal{A}$
is a design parameter: larger symmetry groups produce larger fibers,
smaller $d$, and better generalization per parameter. Training procedures
that reduce curvature---normalization, equivariant architectures,
residual connections---improve not just optimization stability but
compression capacity and generalization simultaneously, because all
three are governed by the same geometric object.

The central open problem is computational: the Jacobian $D\Phi(\theta)$
has dimensions $|\mathcal{N}| \times n$ and is prohibitively large for
modern architectures. Efficient approximations remain the primary
obstacle between this theoretical framework and practice.

What the framework provides is a correct starting point: a precise
identification of the objects that matter and a rigorous account of
why they matter. The geometry of what a neural network learns is the
geometry of the semantic manifold $\mathcal{N}$, not the geometry of
the parameter space $\Theta$. The realization map $\Phi : \Theta \to
\mathcal{N}$ is the structure connecting the two.

\begin{center}
\itshape
Parameters are coordinates in a gauge space. Models are points in the
quotient. Every question about learning, generalization, compression,
and interpretation that is asked in $\Theta$ should be asked instead
in $\mathcal{N}$. Meaning is invariant under implementation.
\end{center}

%% ====================================================================
\bibliographystyle{abbrvnat}
\bibliography{realization_map}

\end{document}
